\documentclass[../../../include/open-logic-section]{subfiles}

\begin{document}

\olfileid[tr]{sth}{ordinals}{vn}	
\olsection[Von Neumann'ın Kuruluşu]{Von Neumann'ın Sıra Sayıları Kuruluşu}

\olref[sth][ordinals][iso]{thm:woalwayscomparable} bir düşünce doğurur.
İyi sıralamaları temsil etmeleri için \emph{sıra tipleri} denen belirli
nesneler tanıtabiliriz. $\tuple{A,<}$ iyi sıralamasının sıra tipi için
$\ordtype{A,<}$ yazarak şu iki ilkeyi elde etmeyi umarız:
\begin{align*}
	\ordtype{A, <} = \ordtype{B, \lessdot} &
	\text{ ancak ve ancak } \ordeq{\tuple{A, <}}{\tuple{B, \lessdot}}\\
	\ordtype{A, <} < \ordtype{B, \lessdot}&
	\text{ ancak ve ancak bir }b \in B\text{ için }
	\ordeq{\tuple{A, <}}{\tuple{B_b, \lessdot_b}}.
\end{align*}
Ayrıca doğal sayıları belirli kümeler olarak tanıtabildiğimiz gibi, sıra
tiplerini de \emph{belirli kümeler olarak} tanıtmayı umabiliriz.

Bunu yapmanın en yaygın yolu---ve bizim izleyeceğimiz yaklaşım---bu sıra
tiplerini belirli \emph{kanonik} iyi sıralanmış kümeler aracılığıyla
tanımlamaktır. Bu kanonik kümeleri ilk kez von Neumann tanıtmıştır:

\begin{defn}
$A$ kümesi, $(\forall x\in A)x\subseteq A$ ise \emph{geçişlidir}.
$A$, ancak ve ancak geçişli ve $\in$ ile iyi sıralanmışsa bir
\emph{sıra sayısıdır}.
\end{defn}
\noindent
Bundan sonra sıra sayıları için Yunan harfleri kullanacağız. Tanımdan
hemen şu çıkar: $\alpha$ bir sıra sayısıysa
$\tuple{\alpha,\in_\alpha}$ bir iyi sıralamadır; burada
$\in_\alpha=\Setabs{\tuple{x,y}\in\alpha^2}{x\in y}$'dir. Böylece
gösterimi biraz kötüye kullanarak $\alpha$'nın \emph{kendisinin} bir iyi
sıralama olduğunu söyleyebiliriz.

İlk birkaç sıra sayımız şunlardır:
\[
	\emptyset, \{\emptyset\},
	\{\emptyset, \{\emptyset\}\},
	\{\emptyset, \{\emptyset\}, \{\emptyset, \{\emptyset\}\}\}, \ldots
\]
Bunların, Sonsuzluk Aksiyomumuzda, yani von Neumann'ın $\omega$ tanımında
karşılaştığımız ilk birkaç sayı olduğunu fark edeceksiniz
(\olref[sth][z][infinity-again]{sec} kesimine bakın). Bu rastlantı değildir.
Von Neumann'ın sıra sayıları tanımı doğal sayıları sıra sayıları olarak ele
alır, fakat sonlu ötesi sıra sayılarına da izin verir.

Her zamanki gibi artık şunu sorabiliriz: sıra sayıları bunlar \emph{mıdır}?
Yoksa von Neumann bize yalnızca sıra sayıları \emph{olarak ele
alabileceğimiz} bazı kümeler mi vermiştir? Bu soru hakkındaki tartışmalar,
\olref[sfr][rel][ref]{sec}, \olref[sfr][arith][ref]{sec},
\olref[sfr][infinite][dedekindsproof]{sec} ve
\olref[sth][z][nat]{sec} içinde yürüttüğümüz tartışmalara benzer; bu yüzden
konuyu uzatmayacağız. Bundan sonra ``sıra sayıları'' ifadesini yalnızca
``von Neumann sıra sayıları'' anlamında kullanacağız.

\end{document}
